\documentclass[10pt,aspectratio=169]{beamer}

% --- PACOTES E CONFIGURAÇÕES DE IDIOMA E FONTE ---
\usepackage[utf8]{inputenc}
\usepackage[portuguese]{babel}
\usepackage[T1]{fontenc}
\usepackage{lmodern}
\usepackage{amsmath,amssymb}
\usepackage{graphicx}
\usepackage{booktabs}
\usepackage{colortbl}

% --- TEMA ACADÊMICO SIMPLES E LIMPO ---
\usetheme{Boadilla}
\usecolortheme{seahorse}

% Customização de cores acadêmicas sóbrias
\definecolor{darkblue}{rgb}{0.1, 0.2, 0.4}
\setbeamercolor{titlelike}{parent=structure,fg=darkblue}
\setbeamercolor{frametitle}{fg=darkblue,bg=gray!10}

\setbeamertemplate{navigation symbols}{} % Remove ícones de navegação poluídos
\setbeamertemplate{footline}[frame number] % Mantém numeração de slides limpa

% --- INFORMAÇÕES DO DOCUMENTO ---
\title[Estimação de Parâmetros via Variograma]{Estimação de Parâmetros Baseada no Variograma}
\subtitle{Avaliação de Desempenho via Simulação Monte Carlo}
\author[Estatística Espacial]{Estudo de Estatística Espacial}
\institute[Moraga (2023)]{Baseado em \textit{Spatial Statistics for Data Science: Theory and Practice with R} (Paula Moraga, Cap. 14)}
\date{\today}

\begin{document}

% ------------------------------------------------------------------------------
% SLIDE 1: TITULO
% ------------------------------------------------------------------------------
\begin{frame}
  \titlepage
\end{frame}

% ------------------------------------------------------------------------------
% SLIDE 2: CONTEXTO E OBJETIVO
% ------------------------------------------------------------------------------
\begin{frame}{Contexto \& Objetivo Estratégico}
  \begin{block}{Motivação}
    \begin{itemize}
      \item A \textbf{Krigagem Ordinária} obtém predições não-viesadas de variância mínima cujos pesos dependem diretamente do modelo de covariância espacial.
      \item Na prática, os parâmetros de covariância (\textbf{Nugget}, \textbf{Partial Sill} e \textbf{Alcance/Range}) precisam ser estimados a partir dos dados observados.
    \end{itemize}
  \end{block}

  \vspace{0.3cm}

  \begin{block}{Objetivo Principal}
    Avaliar quão precisos e robustos são os métodos de estimação baseados no variograma empírico (\textbf{Mínimos Quadrados Ponderados - WLS} e \textbf{Ordinários - OLS}), comparando os parâmetros estimados com os \textbf{parâmetros verdadeiros} sob diferentes tamanhos amostrais ($n = 50, 100, 200$).
  \end{block}
\end{frame}

% ------------------------------------------------------------------------------
% SLIDE 3: DESENHO EXPERIMENTAL
% ------------------------------------------------------------------------------
\begin{frame}{Metodologia da Simulação Monte Carlo}
  \begin{columns}[T]
    \begin{column}{0.48\textwidth}
      \begin{block}{Parâmetros Verdadeiros}
        \begin{itemize}
          \item \textbf{Nugget ($\tau^2$):} $0.20$ (Variância microescala)
          \item \textbf{Partial Sill ($\sigma^2$):} $1.00$ (Variância espacial)
          \item \textbf{Alcance ($\phi$):} $20.0$ (Decaimento espacial)
          \item \textbf{Modelo:} Exponencial $\gamma(h) = \tau^2 + \sigma^2(1 - e^{-h/\phi})$
        \end{itemize}
      \end{block}
    \end{column}

    \begin{column}{0.48\textwidth}
      \begin{block}{Configuração da Simulação}
        \begin{itemize}
          \item \textbf{Réplicas Monte Carlo:} $N = 100$ por cenário
          \item \textbf{Tamanhos Amostrais ($n$):} $50, 100, 200$ pontos
          \item \textbf{Domínio Espacial:} $[0, 100] \times [0, 100]$
          \item \textbf{Métodos de Ajuste:}
            \begin{enumerate}
              \item \textbf{WLS}: Pesos de Cressie $w_k = \frac{N(h_k)}{\gamma(h_k)^2}$
              \item \textbf{OLS}: Pesos unitários $w_k = 1$
            \end{enumerate}
        \end{itemize}
      \end{block}
    \end{column}
  \end{columns}
\end{frame}

% ------------------------------------------------------------------------------
% SLIDE 4: EXEMPLO DE VARIOGRAMA E AJUSTES
% ------------------------------------------------------------------------------
\begin{frame}{Exemplo de Ajuste de Variograma Empírico}
  \begin{center}
    \includegraphics[height=0.73\textheight,keepaspectratio]{plots/exemplo_variograma_ajustado.png}
  \end{center}
\end{frame}

% ------------------------------------------------------------------------------
% SLIDE 5: RESULTADOS VISUAIS (BOXPLOT)
% ------------------------------------------------------------------------------
\begin{frame}{Distribuição dos Parâmetros Estimados}
  \begin{center}
    \includegraphics[height=0.73\textheight,keepaspectratio]{plots/boxplot_parametros.png}
  \end{center}
\end{frame}

% ------------------------------------------------------------------------------
% SLIDE 6: TABELA RESUMO DE RESULTADOS
% ------------------------------------------------------------------------------
\begin{frame}{Resultados Numéricos da Simulação}
  \small
  \begin{table}[h]
    \centering
    \caption{Comparação entre Parâmetros Verdadeiros e Estimados ($N = 100$)}
    \begin{tabular}{ccccccc}
      \toprule
      \textbf{Tamanho ($n$)} & \textbf{Método} & \textbf{Parâmetro} & \textbf{Valor Real} & \textbf{Média Est.} & \textbf{Viés Absoluto} & \textbf{RMSE} \\
      \midrule
      $n = 50$ & OLS & Nugget ($\tau^2$) & 0.20 & 0.218 & +0.018 & 0.198 \\
      $n = 50$ & WLS & Nugget ($\tau^2$) & 0.20 & 0.287 & +0.087 & 0.257 \\
      $n = 50$ & OLS & Alcance ($\phi$)   & 20.0 & 43.76 & +23.76 & 54.62 \\
      \midrule
      $n = 100$ & OLS & Nugget ($\tau^2$) & 0.20 & 0.208 & +0.008 & 0.138 \\
      $n = 100$ & WLS & Nugget ($\tau^2$) & 0.20 & 0.237 & +0.037 & 0.145 \\
      $n = 100$ & OLS & Sill ($\sigma^2$) & 1.00 & 1.376 & +0.376 & 1.133 \\
      \midrule
      $n = 200$ & OLS & Nugget ($\tau^2$) & 0.20 & 0.231 & +0.031 & \textbf{0.109} \\
      $n = 200$ & OLS & Sill ($\sigma^2$) & 1.00 & 1.068 & +0.068 & \textbf{0.625} \\
      $n = 200$ & OLS & Alcance ($\phi$)   & 20.0 & 28.32 & +8.32  & \textbf{30.68} \\
      \bottomrule
    \end{tabular}
  \end{table}
\end{frame}

% ------------------------------------------------------------------------------
% SLIDE 7: CONVERGÊNCIA E EVOLUÇÃO DO RMSE
% ------------------------------------------------------------------------------
\begin{frame}{Convergência e Redução do RMSE com Amostra}
  \begin{center}
    \includegraphics[height=0.73\textheight,keepaspectratio]{plots/rmse_tamanho_amostral.png}
  \end{center}
\end{frame}

% ------------------------------------------------------------------------------
% SLIDE 8: RECOMENDAÇÕES GERENCIAIS
% ------------------------------------------------------------------------------
\begin{frame}{Conclusões \& Recomendações Gerenciais}
  \begin{block}{Principais Descobertas}
    \begin{enumerate}
      \item \textbf{Efeito Amostral Crítico:} Para $n \le 50$, o erro na estimação do alcance ($\phi$) é elevado (RMSE $> 50$). Tamanhos amostrais de pelo menos \textbf{$n \ge 100-200$ pontos} são fortemente recomendados.
      \item \textbf{Estabilidade do Nugget ($\tau^2$):} O efeito pepita apresenta excelente acurácia de estimação mesmo em amostras menores (viés $< 0.02$).
      \item \textbf{WLS vs. OLS:} Para variogramas empíricos de amostras pequenas a moderadas, o ajuste OLS mostrou menor variabilidade extrema na cauda do alcance do que WLS.
    \end{enumerate}
  \end{block}

  \vspace{0.2cm}

  \begin{alertblock}{Recomendações Práticas para Krigagem}
    \begin{itemize}
      \item \textbf{Planejamento Amostral:} Garantir densidade amostral adequada com pares a distâncias curtas e intermediárias.
      \item \textbf{Validação Cruzada:} Sempre validar as predições de Krigagem testando variogramas com limites restritos para evitar superexposição ao alcance.
    \end{itemize}
  \end{alertblock}
\end{frame}

\end{document}
